\begin{afterword}
\summaryhead

The author returns to the mysterious answer of the previous post and names its most critical mistake: the young author did not realize that solving a mystery should make it feel less confusing. The young author knew the stories of astrology, alchemy and vitalism, but thought the chosen phenomenon was a new and stranger kind of mystery, when stars, matter and life had been mysteries for thousands of years before science explained them.

History seemed to teach that astrologers, alchemists and vitalists had a character flaw. Only after seeing the mundane structure inside the mystery did the author realize how reasonable vitalism had seemed at the time, and how surprising the universe's answer was. We read history but do not live it. Had the author lived through those discoveries, the post concludes, the new mysterious answer would have been rejected at once.

\responsehead

The post is short, and its lesson is sound as far as it goes: an answer that leaves a phenomenon as mysterious as before has not explained it. Two things limit it.

The rule catches only one kind of error. It flags an answer that still feels mysterious. It does not flag an answer that feels like an explanation and is not, which is the error the post of the day before, ``Say Not Complexity,'' describes. A reader who applies only this rule will catch the young author's kind of mistake and miss the other kind.

The remedy is a counterfactual that the book's own example does not support. The post claims that, had the author lived through the fall of astrology, alchemy and vitalism, the new mysterious answer would have been recognized at once. The vitalist the author quoted five days earlier, Lord Kelvin, did live through a scientific revolution in his own field. He gave up Carnot's theory of heat for Joule's, and helped formulate the laws of thermodynamics. He still held that the influence of life on matter was ``infinitely beyond the range of any scientific inquiry hitherto entered on.'' Living through history at first hand did not protect him.

The history itself is not told. How reasonable vitalism seemed ``at the time'' is asserted, and the reader must rely on the earlier post for the one vitalist named. The claim about how ``the children of that generation'' greet each new puzzle is based only on the author's own case.

In short: a sound warning sign for one kind of empty answer, together with a remedy, living through the history of science, that the case of Kelvin, the vitalist the book quotes, counts against.
\end{afterword}
